% !TEX root = ../../mixed-sklyanin-cubics.tex
% Current result: 2026-11-23 / v69.
\subsection{階数一の中心 norm の全特異判別と軸外の Azumaya 二重点}
\label{ifn:subsec:full-norm}

階数一の係数スライスでは、前小節の軸上の四次判別に加えて
横二変数組の非零臨界枝がある。この四種類の枝を全て用いれば、
中心 norm が定義する \emph{可換}三次超曲面の特異条件を
必要十分条件として記述できる。この結果を残余圏の全 smooth
判別との同一視には用いない。

命題\ref{ifd:prop:ambient-locus}の記号を用い、
$\Lambda=\lambda E_{11},\,M=\mu E_{11}$、
$t=\tau_1,\,u=\tau_2$、$a=\alpha_1^2,\,b=\alpha_2^2$ とする。
$t,u,a,b$ は非零であり、$X=X_1$、$Y=Y_1$ と略記する。
四つの線型式と二変数三次式を
\begin{equation}\label{ifn:eq:branch-functions}
 \begin{split}
 P(X,Y)&=aX(\lambda Y-tX)^2-bY(\mu X-uY)^2,\\
 k_x(X,Y)&=\frac{(4-t^3)X+t^2\lambda Y}{3t^2},\qquad
 k_y(X,Y)=\frac{(4-u^3)Y+u^2\mu X}{3u^2},\\
 F_{\epsilon,\delta}(X,Y)
   &=P(X,Y)-\epsilon\,at^2k_x(X,Y)^3
               +\delta\,bu^2k_y(X,Y)^3,
 \qquad (\epsilon,\delta)\in\{0,1\}^2
 \end{split}
\end{equation}
と置く。$F_{\epsilon,\delta}$ は斉次三次式である。
$\operatorname{Disc}(F)$ は binary cubic の通常の判別式とする。
ここではその零集合だけを扱い、判別式の重複度は主張しない。

\begin{theorem}[中心 norm 超曲面の全特異判別]
\label{ifn:thm:four-discriminants}
六平方座標の射影空間 $\mathbb P^5$ 内の
\[
  H_{\lambda,\mu}=V(\mathcal N)\subset\mathbb P^5,
  \qquad
  \mathcal N=-a\det q_x+b\det q_y
\]
について、複素係数の点ごとに
\begin{equation}\label{ifn:eq:exact-discriminant}
 H_{\lambda,\mu}\ \text{が特異}
 \quad\Longleftrightarrow\quad
 \prod_{\epsilon,\delta\in\{0,1\}}
   \operatorname{Disc}(F_{\epsilon,\delta})=0
\end{equation}
が成立する。

より正確に、特異点の横座標は全て
\begin{equation}\label{ifn:eq:critical-branches}
 (X_2,X_3)=
 \begin{cases}
 (0,0),&\epsilon=0,\\
 (k_x\xi,k_x\xi^2),\ \xi^3=1,&\epsilon=1,
 \end{cases}
 \qquad
 (Y_2,Y_3)=
 \begin{cases}
 (0,0),&\delta=0,\\
 (k_y\eta,k_y\eta^2),\ \eta^3=1,&\delta=1
 \end{cases}
\end{equation}
であり、対応する $[X:Y]\in\mathbb P^1$ は
$F_{\epsilon,\delta}$ の重根である。
$k_x=0$ または $k_y=0$ の場合には二つの枝が合流してもよい。
この必要十分条件は $\mathcal N$ が定義する中心 hypersurface の
特異性についてのもので、中心の非 Azumaya 層の構造、
非可換残余圏の全 smooth locus または可換 Ku 標的の判別を
決めたという主張ではない。
\end{theorem}

\begin{proof}
式\eqref{ifb:eq:clifford-matrices}の二つの行列式を展開すると、
六変数三次式は厳密に
\begin{equation}\label{ifn:eq:complete-norm}
 \begin{split}
 4\mathcal N={}&P(X,Y)
 +at^2\bigl(X_2^3+X_3^3-3k_xX_2X_3\bigr)\\
 &-bu^2\bigl(Y_2^3+Y_3^3-3k_yY_2Y_3\bigr)
 \end{split}
\end{equation}
となる。横偏微分を零に置くと
\[
 X_2^2=k_xX_3,\quad X_3^2=k_xX_2,\qquad
 Y_2^2=k_yY_3,\quad Y_3^2=k_yY_2
\]
である。$t,u,a,b$ は非零なので、解は式
\eqref{ifn:eq:critical-branches}に尽きる。
特に $X=Y=0$ なら $k_x=k_y=0$ であり、
全横座標も零となるから、射影的な臨界点は存在しない。

式\eqref{ifn:eq:critical-branches}のいずれかの枝に
$4\mathcal N$ を代入すると
\[
 4\mathcal N|_{\epsilon,\delta}
 =P-\epsilon at^2k_x^3+\delta bu^2k_y^3
 =F_{\epsilon,\delta}
\]
となる。
この枝に沿う横偏微分はすでに零なので、chain rule により
$\partial_X(4\mathcal N)|_{\epsilon,\delta}
 =\partial_XF_{\epsilon,\delta}$ および
$\partial_Y(4\mathcal N)|_{\epsilon,\delta}
 =\partial_YF_{\epsilon,\delta}$
である。従って六変数の全偏微分が零となる条件は、
ある枝で $\partial_XF_{\epsilon,\delta}
 =\partial_YF_{\epsilon,\delta}=0$ となる条件と等価である。
斉次三次式の Euler 等式により、これらが零なら
$F_{\epsilon,\delta}=0$ でもある。
標数零の二変数斉次三次式について、これは
$\mathbb P^1$ に重根があること、
すなわち判別式が零であることと同値である。
$\xi,\eta$ の三つの選択は $k_x^3,k_y^3$ と横一次微分に
同じ値を与えるため、判別式の積の零集合が
全六変数中心 hypersurface の特異係数集合と一致する。
\end{proof}

最初の枝 $(\epsilon,\delta)=(0,0)$ だけを取ると
$F_{0,0}=P$ であり、前小節の
式\eqref{ifq:eq:disc-factor}を回収する。
残る三枝は軸外の特異点を検出する。

\begin{corollary}[軸上四次曲線で検出できない Azumaya 特異点]
\label{ifn:cor:off-axis-azumaya}
$\tau_1=\tau_2=2$、$a_1=a_2=1$、$\lambda=0$ とし、
$\mu\in\mathbb C$ は $\mu^3=-196$ の任意の根とする。
周囲代数は $h=49\ne0$ のため六次元 AS 正則であり、
中心三次の普遍正則係数は $c=196\ne0$ である。
一方、$H_{0,\mu}$ には
\begin{equation}\label{ifn:eq:off-axis-point}
 [X_1:X_2:X_3:Y_1:Y_2:Y_3]
 =
 [28/(13\mu):0:0:1:5/13:5/13]
\end{equation}
という通常二重点がある。この点では
$\det q_x=\det q_y=-(28/(13\mu))^3\ne0$ であるから、
射影中心の Azumaya 開に位置する。
この点において、規格化した $4\mathcal N/196$ の
$Y_1=1$ chart での五変数 Hessian 行列式は
\begin{equation}\label{ifn:eq:off-axis-hessian}
                 \det(\operatorname{Hess})=\frac{12902400}{2197}\ne0
\end{equation}
である。前小節の軸上四次因子 $Q$ は
$Q/\alpha_1^4=-2000\ne0$ なので、
この点は前小節の $V(Q)$ では検出されない。

従って、前小節の $V(Q)$ に含まれない
AS 正則かつ正則中心三次の係数でも、
残余線束圏が非 smooth になる。
\end{corollary}

\begin{proof}
両因子の規格化 $\alpha_1^2=\alpha_2^2=196$ を
共通因子として除く。$\lambda=0$、
$\mu^3=-196$ では $h=49$ と $c=196$ は
命題\ref{ifd:prop:ambient-locus}から従う。
$Y=1$ とおき $z=\mu X$ とすると
\[
 \frac{F_{0,1}(X,1)}{196}
 =\frac{(z-7)(13z-28)^2}{1323}.
\]
従って $X=28/(13\mu)$ は二重根で、
$k_y=(z-1)/3=5/13$ であり、$X_2=X_3=0$ は第零枝を選んだものである。
定理\ref{ifn:thm:four-discriminants}により
式\eqref{ifn:eq:off-axis-point}が特異点となる。
両行列式は表示した同じ非零数であり、
Hessian は式\eqref{ifn:eq:complete-norm}から直接計算すると
式\eqref{ifn:eq:off-axis-hessian}となる。

$Q$ は前小節の式\eqref{ifq:eq:quartic}であり、
$\alpha_1^2=\alpha_2^2=196$ を代入して共通因子
$\alpha_1^4$ で割れば
$Q/\alpha_1^4=8\mu^3-432=-2000$ となる。
非退化Clifford点の射影Azumaya chart では、
完備局所代数がこの通常二重点の完備中心環上の
四次行列代数となる。前小節の
定理\ref{ifq:thm:axis-discriminant}と同じ
無限自己Ext、coherent 非 proper と
smooth 線束残余の飽和矛盾が適用できる。
\end{proof}


% !TEX root = ../../mixed-sklyanin-cubics.tex
% Current result: 2026-11-23 / v69.
\subsubsection{任意の二因子に対する一方向階数一の特異値と局所型}
\label{ifod:subsubsec:oneway}

前定理の対称特殊化を任意の許容 Clifford 二因子へ拡張する。
係数を一方向の $\Lambda=sE_{11},M=0$ に制限すると、
周囲の AS 正則性と中心元の正則性は全 $s$ で保たれる。
一方、中心 norm 超曲面の特異係数は、四つの三次方程式で正確に求まる。
本小節の判別は可換中心 norm の判別であり、固定 Hesse 切片の
高次 Ku 比較係数と同一視しない。

命題\ref{ifd:prop:ambient-locus}と定理\ref{ifn:thm:four-discriminants}の記法を用い、
\[
 \Lambda=sE_{11},\quad M=0,\qquad
 t=\tau_1,\quad u=\tau_2,\quad
 a=\alpha_1^2,\quad b=\alpha_2^2
\]
と置く。$t,u\ne0$、$t^3,u^3\notin\{1,-8\}$、$a,b\ne0$ とする。
\begin{equation}\label{ifod:eq:g}
 g=\frac{(u^3-1)(u^3+8)^2}{27u^6}
   =1-\frac{(4-u^3)^3}{27u^6},\qquad
 g_0=1,\quad g_1=g.
\end{equation}
従って $g\ne0$ である。四つの非零数を
\begin{equation}\label{ifod:eq:critical-cubes}
 \begin{aligned}
 \rho_{0\delta}&=\frac{27btu^2}{4a}g_\delta,\\
 \rho_{1\delta}&=-\frac{bu^2(t^3-1)^2(t^3+8)}{4at^2}g_\delta,
       \qquad \delta\in\{0,1\}
 \end{aligned}
\end{equation}
で定める。

\begin{theorem}[一方向階数一の全中心 norm 特異係数]
\label{ifod:thm:general-values}
上の任意の許容因子と非零重みに対し、中心 norm の射影超曲面
$H_s=V(\mathcal N_s)\subset\mathbb P^5$ は
\begin{equation}\label{ifod:eq:general-discriminant-set}
 H_s\text{ が特異}\quad\Longleftrightarrow\quad
       \prod_{\epsilon,\delta\in\{0,1\}}
          (s^3-\rho_{\epsilon\delta})=0
\end{equation}
を満たす。全 $s$ で $h=(1-t^3)(1-u^3)\ne0$、
$c=at^2/4\ne0$ なので、周囲代数は六次元 Koszul AS 正則・非捩れ
Calabi--Yau であり、指定中心三次も普遍的に正則である。

式\eqref{ifod:eq:general-discriminant-set}の全ての特異点は
$Y_1\ne0$ の chart 上にあり、両 Clifford 行列式が非零の
Azumaya 開に属する。従ってその係数における
全 coherent 残余圏は非 proper、残余線束圏は非 smooth である。
この逆向き、すなわち $H_s$ が滑らかなときの残余圏の
smoothness は本定理では主張しない。
\end{theorem}

\begin{proof}
$Y_1=1$ とし、$X=X_1$ と書く。
四臨界枝\eqref{ifn:eq:critical-branches}のうち
$\epsilon$ が第一因子、$\delta$ が第二因子の非零横臨界枝を表す。
$M=0$ では
\[
 k_y=\frac{4-u^3}{3u^2},\qquad
 k_x=\frac{(4-t^3)X+t^2s}{3t^2}.
\]
第二因子の横臨界値は $-bu^2g_\delta$ なので、
式\eqref{ifn:eq:branch-functions}の四つの二変数三次は
\begin{equation}\label{ifod:eq:reduced-cubics}
 F_{\epsilon\delta}(X,1)=a\phi_\epsilon(X)-bu^2g_\delta,
\quad
 \phi_0=X(s-tX)^2,\qquad
 \phi_1=\phi_0-\frac{((4-t^3)X+t^2s)^3}{27t^4}.
\end{equation}
$F_{0\delta}$ の $X^3$ 係数は $at^2$、
$F_{1\delta}$ のそれは
$at^2(t^3-1)(t^3+8)^2/(27t^6)$ で、いずれも非零である。
従って全ての射影特異点は $Y_1\ne0$ であり、
二変数三次の判別に有限 chart の三次多項式を使える。

$\phi_0'=(s-tX)(s-3tX)$ であり、値が零でない臨界点と
その臨界値は
\[
 X_0=\frac{s}{3t},\qquad
 \phi_0(X_0)=\frac{4s^3}{27t}.
\]
さらに直接微分すると
\begin{equation}\label{ifod:eq:phi-one-derivative}
 \phi_1'(X)=
 \frac{((t^3+8)X-t^2s)
   ((t^3+8)(t^3-1)X-t^2(t^3+5)s)}{9t^4}.
\end{equation}
第一の臨界点 $X=t^2s/(t^3+8)$ での値は零である。
値が零でない第二の臨界点とその値は
\begin{equation}\label{ifod:eq:second-value}
 X_1=\frac{st^2(t^3+5)}{(t^3-1)(t^3+8)},\qquad
 \phi_1(X_1)=-\frac{4s^3t^2}{(t^3-1)^2(t^3+8)}.
\end{equation}
$b u^2g_\delta/a\ne0$ なので、式\eqref{ifod:eq:reduced-cubics}の
重根は $\phi_\epsilon(X_\epsilon)=bu^2g_\delta/a$ と同値である。
これは正確に式\eqref{ifod:eq:critical-cubes}となる。
別法として、二変数三次の標準判別式を展開すれば
\begin{equation}\label{ifod:eq:factor-four-discriminants}
 \begin{aligned}
 \operatorname{Disc}(F_{0\delta})
   &=4a^3bt^3u^2g_\delta(s^3-\rho_{0\delta}),\\
 \operatorname{Disc}(F_{1\delta})
   &=-\frac{4a^3bu^2g_\delta(t^3+8)^3}{27t^6}
                       (s^3-\rho_{1\delta})
 \end{aligned}
\end{equation}
を得る。前定理\ref{ifn:thm:four-discriminants}を適用すれば
必要十分条件\eqref{ifod:eq:general-discriminant-set}が従う。
$h,c$ の主張は命題\ref{ifd:prop:ambient-locus}による。

Azumaya 性も全係数に対して直接確認できる。
$Y_1=1$ の第二因子の臨界枝で、
\begin{equation}\label{ifod:eq:nonzero-determinants}
              \det q_y=-\frac{u^2g_\delta}{4}\ne0.
\end{equation}
$\mathcal N=-a\det q_x+b\det q_y=0$ だから
$\det q_x=(b/a)\det q_y\ne0$ である。
従って両 Clifford 行列は非退化であり、
定理\ref{ifq:thm:axis-discriminant}の射影Azumaya chartを適用できる。
中心 norm の特異局所環の剰余体は有限射影次元を持たず、
その点対象の無限自己Ext、全 coherent 残余非proper、
有限線束残余の非smooth性は同定理の証明と同じである。
\end{proof}

\begin{theorem}[一般因子での節点数と判別式の重複度]
\label{ifod:thm:generic-nodes}
さらに
\begin{equation}\label{ifod:eq:node-open}
                t^3\ne-2,\qquad u^3\ne4
\end{equation}
を仮定する。このとき $s^3=\rho_{\epsilon\delta}$ に現れる特異点は
すべて通常二重点である。
四つの $\rho_{\epsilon\delta}$ が相異なれば、
特異係数はちょうど十二個で、各ファイバーの節点数は
\[
             3^{\epsilon+\delta}=1,3,3,9
\]
となる。異なる $\rho$ が一致する場合でも、異なる臨界枝の点は
合流せず、同じファイバー内の節点数を足し合わせればよい。

六変数三次超曲面の通常の判別多項式を
$s$ に制限したものは、非零 scalar を除き
\begin{equation}\label{ifod:eq:scheme-disc}
 \boxed{\quad
 \mathfrak D_{t,u,a,b}(s)
   \doteq \prod_{\epsilon,\delta\in\{0,1\}}
            (s^3-\rho_{\epsilon\delta})^{3^{\epsilon+\delta}}
 \quad}
\end{equation}
である。ここで $\doteq$ は $\C^\times$ の倍を除いた等号である。
この多項式の $s$ に関する次数は $48$ であり、
零集合だけでなく通常の hypersurface discriminant の
スキーム的重複度を記述する。
\end{theorem}

\begin{proof}
第一因子の二つの非零臨界値に対応する $X$ の二階微分は
\[
 \phi_0''(X_0)=-2ts,\qquad
 \phi_1''(X_1)=\frac{2s(t^3+8)}{3t^2},
\]
であり非零である。さらに横臨界枝での $k_x$ の値は
\begin{equation}\label{ifod:eq:transverse-k}
 k_x(X_0)=\frac{2s(t^3+2)}{9t^3},\qquad
 k_x(X_1)=\frac{2s(t^3+2)}{(t^3-1)(t^3+8)},\qquad
 k_y=\frac{4-u^3}{3u^2}.
\end{equation}
式\eqref{ifod:eq:node-open}では、これらは全て非零である。
横二変数組 $z_2^3+z_3^3-3kz_2z_3$ の Hessian は
原点の枝で行列式 $-9k^2$、非零枝 $(k\xi,k\xi^2)$ で
$27k^2$ を持つ。横変数の臨界切断に沿う Schur complement は
$a\phi_\epsilon''(X_\epsilon)$ なので、
五変数の全 Hessian は非退化である。
$\epsilon=1$ には三つの独立した $\xi^3=1$ の選択、
$\delta=1$ には三つの独立した $\eta^3=1$ の選択がある。
式\eqref{ifod:eq:node-open}では非零枝は零枝と一致しない。
よって節点数は $3^{\epsilon+\delta}$ となる。

各節点が $s$ に関して一次で平滑化されることも必要である。
全六変数の式を微分し、臨界切断に制限すれば
\begin{equation}\label{ifod:eq:parameter-transversality}
 \left.\partial_s(4\mathcal N)\right|_{\epsilon=0}
      =\frac{4as^2}{9t}\ne0,\qquad
 \left.\partial_s(4\mathcal N)\right|_{\epsilon=1}
      =-\frac{12as^2t^2}{(t^3-1)^2(t^3+8)}\ne0.
\end{equation}
四つの $\rho$ は非零なので $s\ne0$ である。
標数零の形式 Morse 補題により、各節点の一パラメータ完備局所族は
$\sum_{i=1}^5z_i^2+v(s)=0$ と書け、
$v'(s_0)\ne0$ である。従って通常の hypersurface discriminant の
局所方程式は各節点で $v(s)$ の一次因子を一つ持つ。
同じファイバーに複数の独立な節点があれば、その局所因子の積である。
定理\ref{ifod:thm:general-values}は他の特異係数がないことを示すので、
$s$ に関する非零多項式としての判別式は表示した積の定数倍に等しい。
次数は $3(1+3+3+9)=48$ である。
\end{proof}

\begin{proposition}[許容範囲内の枝の衝突と $D_4$ 特異点]
\label{ifod:prop:collision-types}
式\eqref{ifod:eq:critical-cubes}の二つの行が一致する条件は
$t^3=-2$、二つの列が一致する条件は $u^3=4$ である。
これらはいずれも許容される Clifford 因子の特殊値である。

$t^3=-2$、$u^3\ne4$ なら、各 $s^3=\rho_{0\delta}=\rho_{1\delta}$ の
Azumaya 特異点は、$\delta=0$ では一つ、$\delta=1$ では三つで、
それぞれ非退化二次形式を三変数分と
$x_2^3+x_3^3$ を持つ $D_4$ 型である（Milnor 数 $4$）。
$u^3=4$、$t^3\ne-2$ なら、各 $s^3=\rho_{\epsilon0}=\rho_{\epsilon1}$ の
Azumaya 特異点は $\epsilon=0$ で一つ、$\epsilon=1$ で三つであり、
同様に $D_4$ 型である。
両条件が成立すると四つの $\rho$ は全て一致し、
各 $s^3=\rho$ に一つの Azumaya 特異点があり、
その局所関数は非退化一変数二次式と四変数の
$x_2^3+x_3^3+y_2^3+y_3^3$ の和に形式的に同値で、
Milnor 数は $16$ である。
この命題は中心normの局所型を述べるものであり、
可換 Ku 比較の特殊化や全 residual smoothness を主張しない。
\end{proposition}

\begin{proof}
$t^3=-2$ では恒等式
\[
 27t^3+(t^3-1)^2(t^3+8)=(t^3+2)^3
\]
により二つの非零臨界値は一致し、式\eqref{ifod:eq:transverse-k}の
二つの $k_x$ は零になる。このとき $X_0=X_1=s/(3t)$ であり、
$X=s/(3t)+v$ とおくと $Y_1=1$ chart の第一因子の関数は
\[
 a\left(\mathrm{const}-tsv^2+t^2v^3
       +t^2(x_2^3+x_3^3)-6vx_2x_3\right).
\]
$v$ の二次係数は非零なので、形式 Morse 消去後の
$x_2,x_3$ の最低次は非退化な三次 $x_2^3+x_3^3$ である。
四次以上の補正はその偏微分 $(x_2^2,x_3^2)$ を用いる
次数ごとの座標変更で除去できる。
これは $D_4$ 型で、Jacobian 局所環の長さは $4$。
$u^3\ne4$ なら第二因子の横二変数は Morse 型なので、
$\delta=0$ で一点、$\delta=1$ で三点となる。

$u^3=4$ では $k_y=0$ かつ $g=1$ で、
第二因子の横関数は正確に $-bu^2(y_2^3+y_3^3)$ となる。
$t^3\ne-2$ では第一因子は非退化な三変数二次形式へ消去できる。
従って $\epsilon=0$ の一つ、$\epsilon=1$ の三つが
同様の $D_4$ 型となる。

両方の条件が成り立てば、二つの二変数三次型が残り、
一次変数 $v$ を消去する際の四次補正は
$(x_2^2x_3^2)$ に属する。この補正は上記の三次の偏微分
$(x_2^2,x_3^2,y_2^2,y_3^2)$ を使って除去できる。
従って表示した四変数 Fermat 三次型の suspension となり、
Jacobian 局所環の長さは $2^4=16$ である。
\end{proof}

\begin{remark}[規約と高次比較の残存問題]
\label{ifod:rem:scope}
$\tau_1^3=-2$ は標準 Hesse の Fermat 因子であり、
$\tau_2^3=4$ は中心 norm の第二横Hesse関数で二次結合係数が
消える特殊点である。この二つを混同しない。
また二組の $\rho$ が交差して一致することは一般にもあり得るが、
式\eqref{ifod:eq:node-open}の下では節点は別々の点に留まる。
中心 norm の特異係数・scheme 的重複度・Milnor 型の記述から
原点の固定 Hesse 切片の十二の高次比較係数は決定できない。
必要なのは指定した実核・Picard・周期を含む高次変形の移送である。
本小節の恒等式と一般例の Hessian は
\texttt{computations/verify\_clifford\_general\_oneway.py}
で厳密に検算できる。
\end{remark}


\begin{theorem}[対称係数の四判別式と一変数probeの十二特異値]
\label{ifn:thm:symmetric-probe}
$\tau_1=\tau_2=2$、$a_1=a_2=1$ とする。
$\alpha_1^2=\alpha_2^2=196$ を除いた
二変数三次 $\widehat F_{\epsilon,\delta}=F_{\epsilon,\delta}/196$ とその判別式
$D_{\epsilon,\delta}=\operatorname{Disc}_X\widehat F_{\epsilon,\delta}(X,1)$ を考える。
$q=\lambda\mu$ と置くと、
\begin{equation}\label{ifn:eq:symmetric-discriminants}
\begin{aligned}
 D_{00}&=(q-4)^2Q_{00},&
 D_{01}&=-\frac{(q-8)^2}{27}Q_{01},\\
 D_{10}&=-\frac{(q-8)^2}{27}Q_{10},&
 D_{11}&=-\frac{(q-16)^2}{19683}Q_{11},
\end{aligned}
\end{equation}
ただし
\begin{equation}\label{ifn:eq:symmetric-residuals}
\begin{aligned}
 Q_{00}={}&q^2+8(\lambda^3+\mu^3)+72q-432,\\
 Q_{01}={}&16\lambda^4\mu-56\lambda^3+157q^2-1456q
                   +16\mu^3+3136,\\
 Q_{10}={}&16\lambda^3+157q^2+16\lambda\mu^4
                   -1456q-56\mu^3+3136,\\
 Q_{11}={}&256q^4-8032q^3+87117q^2-393568q+614656\\
          &+(3024-432q)(\lambda^3+\mu^3)
\end{aligned}
\end{equation}
である。ここでは各 $Q_{\epsilon,\delta}$ の絶対既約性は主張しない。

二乗因子 $q=4,8,16$ は、二つの Clifford 行列式がともに零となる
特異中心点を持つ。これらの因子と別に、四つの
$Q_{\epsilon,\delta}$ のそれぞれに、Azumaya 開での通常二重点を
持つ係数が存在する。

特に一方向の行列単位probe
\begin{equation}\label{ifn:eq:probe-line}
                 \Lambda=sE_{11},\qquad M=0
\end{equation}
上では、周囲代数は全 $s\in\C$ で六次元 AS 正則かつ非捩れ
Calabi--Yau であり、指定中心三次は全 $s$ で正則である。
中心 norm 超曲面 $H_s=V(\mathcal N_s)\subset\mathbb P^5$ は
\begin{equation}\label{ifn:eq:probe-bad}
 (s^3-54)(s^3-56)(s^3+196)(27s^3+5488)\ne0
\end{equation}
のとき、かつそのときに限り滑らかである。
零となる十二の係数点では、$H_s$ の特異点は全て
Azumaya 開の通常二重点であり、各fiberでの個数は
\begin{equation}\label{ifn:eq:probe-node-counts}
\begin{array}{c|cccc}
 s^3&54&56&-196&-5488/27\\\hline
 \text{通常二重点数}&1&3&3&9
\end{array}
\end{equation}
である。
したがって、これら十二の係数では全 coherent 残余圏は非 proper、
残余線束圏は非 smooth となる。
条件\eqref{ifn:eq:probe-bad}で残余圏が smooth とまでは主張しない。
\end{theorem}

\begin{proof}
定理\ref{ifod:thm:general-values}の特殊化を含む、式\eqref{ifn:eq:branch-functions}の $t=u=2$ と $a=b=1$ の特殊化を
$\widehat F_{\epsilon,\delta}$ として、binary cubic の判別式
\[
  \operatorname{Disc}(aX^3+bX^2+cX+d)
  =b^2c^2-4ac^3-4b^3d-27a^2d^2+18abcd
\]
を展開すれば、式\eqref{ifn:eq:symmetric-discriminants}と
\eqref{ifn:eq:symmetric-residuals}が得られる。

二乗因子の幾何学的意味も直接確認できる。
$Y_1=1$、$\lambda\ne0$ の chart において
四臨界枝に次の点を取る：
\[
\begin{array}{c|c|c|c}
(\epsilon,\delta)&\lambda\mu&X_1&
 (X_2,X_3;\,Y_2,Y_3)\\\hline
(0,0)&4&\lambda/2&(0,0;\,0,0)\\
(0,1)&8&\lambda/2&(0,0;\,1,1)\\
(1,0)&8&\lambda/4&(\lambda/4,\lambda/4;\,0,0)\\
(1,1)&16&\lambda/4&(\lambda/4,\lambda/4;\,1,1)
\end{array}
\]
これらを式\eqref{ifb:eq:clifford-matrices}に代入すると、
$\mathcal N$ と全偏微分が零であり、
$\det q_x=\det q_y=0$ である。従って表示した二乗因子は
非 Azumaya の臨界枝を持つ。

四つの $Q_{\epsilon,\delta}$ に現れる独立の Azumaya
特異点を示すため、逆向きの一方向probe
$\lambda=0,\ \mu=m$ に制限する。
式\eqref{ifn:eq:symmetric-discriminants}は
\begin{equation}\label{ifn:eq:probe-four-discriminants}
\begin{aligned}
 D_{00}&=128(m^3-54),&
 D_{01}&=-\frac{1024}{27}(m^3+196),\\
 D_{10}&=\frac{3584}{27}(m^3-56),&
 D_{11}&=-\frac{28672}{19683}(27m^3+5488)
\end{aligned}
\end{equation}
となる。四つの三乗値は相異なり、どれも零でない。
$Y_1=1$、$z=mX_1$ と置くと、各場合の binary cubic は
ただ一つの二重根を持つ。重根と、横座標の
位相 $(\xi,\eta)=(1,1)$ での
Clifford 行列式、規格化した $4\mathcal N/196$ の
五変数 Hessian 行列式は次の通りである：
\begin{equation}\label{ifn:eq:probe-witnesses}
\begin{array}{c|c|c|c|c}
(\epsilon,\delta)&m^3&z&
 \det q_x=\det q_y&\det\operatorname{Hess}\\\hline
(0,0)&54&6&-4&153600\\
(0,1)&-196&28/13&112/2197&12902400/2197\\
(1,0)&56&6&-4&-460800\\
(1,1)&-5488/27&28/13&112/2197&-38707200/2197
\end{array}
\end{equation}
横座標は式\eqref{ifn:eq:critical-branches}で、
$k_x=-z/(3m)$、$k_y=(z-1)/3$ とすることにより復元できる。
表の値は行列式を展開し、$m^3$ を各値に置き換えた
有理数恒等式である。Hessian は全て非零であり、二つの
Clifford 行列式も非零なので、これらは Azumaya 開の
通常二重点である。

$\epsilon=1$ の場合の $\xi^3=1$ の三択と、
$\delta=1$ の場合の $\eta^3=1$ の三択は、
それぞれ別々の射影点を与える。
これらの座標変換は $\mathcal N$ を保ち、Jacobian の行列式が一
なので、Hessian の非退化性も保持する。
したがって各特異fiberの点数は $1,3,3,9$ となる。
他の臨界枝の判別式はその係数では零とならず、
重根も一つだけなので、これ以外に特異点はない。

$\lambda=s$、$\mu=0$ の族は、二因子を交換すれば
$\lambda=0$、$\mu=s$ の族に一致する。
中心重みも等しいため、二つの norm の全体の符号が変わるだけで
零集合・局所型・点数は一致する。
また命題\ref{ifd:prop:ambient-locus}により
$h=49$、$c=196$ は全probe上で非零である。
これで中心normの滑らかさの必要十分条件と点数を得る。
各通常二重点が Azumaya chart にあることから、
系\ref{ifn:cor:off-axis-azumaya}および
定理\ref{ifq:thm:axis-discriminant}と同じ点対象・coherent
飽和の論法により残余圏の非smooth性が従う。
\end{proof}

\begin{remark}[一変数probeと高次比較係数]
\label{ifn:rem:probe-vs-jets}
式\eqref{ifn:eq:probe-line}は
定理\ref{ifj:thm:twelve-coefficients}の七つの
probe のうちの行列単位方向である。
本定理はその中心normの特異係数と各fiberの局所型を
\emph{全係数}で決定するが、分離原点の固定した
Hesse多項式切片への実Ku比較の二次・三次Taylor係数を
決定していない。原点の形式変形空間では
判別因子が原点を通らないので、その方程式だけで
未計算の十二係数を復元することはできない。
原点の比較核、Picard補正と指定周期を保つ
実Bockstein移送の高次計算が必要である。
本定理の多項式・Hessian・位相点数は
\texttt{computations/verify\_clifford\_symmetric\_discriminants.py}
で再現できる。
\end{remark}

\begin{remark}[判別因子の数と残る境界]
\label{ifn:rem:scope}
$t=u=2$、$a=b=1$ の規格化で、第二の枝の判別式は
\[
 \begin{split}
 \operatorname{Disc}(F_{0,1})=
 -\frac{(\lambda\mu-8)^2}{27}
 \bigl(&16\lambda^4\mu-56\lambda^3
 +157\lambda^2\mu^2-1456\lambda\mu\\
 &+16\mu^3+3136\bigr)
 \end{split}
\]
である。$\lambda=0,\mu^3=-196$ は後者の因子上にあり、
$\lambda\mu=8$ という非 Azumaya 二重因子上にはない。
四つの binary cubic の判別式を用いれば、
中心 norm の射影 hypersurface の全特異係数集合を記述できるが、
この四つの因子の全てについて
非 Azumaya 中心の実圏模型を構成したわけではない。
また、残余smoothnessが中心normの正則性と同値であること、
全係数の可換 Ku 比較、固定Hesse切片の十二の高次係数の値は
未証明・未計算のままである。

\texttt{computations/verify\_clifford\_fullnorm.py} は、
式\eqref{ifn:eq:complete-norm}の六変数恒等式、
四臨界枝の偏微分、四判別式の定義、
新しい Azumaya 二重点の値と Hessian を
有理数・代数的数の厳密計算として検査する。
有限計算だけで残余圏の非 smooth 性を証明したとはしない。
\end{remark}
